\subsection{Sampling, inputs and extraction}
The initial BigCodeBench pool contains 1,140 tasks. Task IDs are sorted as strings and shuffled with Python's \texttt{random.Random(20260908)}. The first 40 form the development sample. From the remaining order, the initial factorial study takes the first 200 IDs after excluding /0--/7; the independent layout study takes the next 80 after the same exclusions. The complete outcome tables identify these samples. In the initial factorial study, each of four 50-task shards is sorted by task ID before task blocks and the five conditions within a block are shuffled with the same ordering seed. Repetition labels distinguish fresh generations and do not control the model's random sampling.

For BigCodeBench, the task block contains the original \texttt{instruct\_prompt}; the supplied-context block contains the original \texttt{code\_prompt}, including starter code. Neither the \texttt{test} field nor the \texttt{canonical\_solution} enters the model input. The main prompt appendix specifies the operation sentences, labels, block order and final-output instruction. In the layout experiment, the displayed treatment instruction follows the task text and precedes the supplied-context block.

For D, SN, SR, MN, MR and INoT*, code extraction recognizes fenced blocks with the case-insensitive tag \texttt{python} or \texttt{py}, a newline after the opening fence and a newline before the closing fence. Exactly one recognized block is required; other language tags are ignored by the extractor. Leading and trailing whitespace is stripped from that block. Zero or multiple recognized blocks yield an empty candidate and a format-error flag. This implementation is narrower than interpreting or repairing arbitrary model formatting. Empty candidates are evaluated and retained as failures when a valid completed response exists; absent generations remain unavailable. SCC uses its author's extraction procedure described below.

\subsection{Native tests and environment}
The initial factorial evaluator uses Python 3.11, BigCodeBench v0.1.4 tasks, local execution, \texttt{instruct full}, two parallel evaluator workers and the selected task IDs. Automatic insertion of starter code is disabled (\texttt{calibrated=False}); separately completed reference controls replace automatic reference generation (\texttt{no\_gt=True}). The nominal minimum-time option is 0.1 seconds, but the recorded effective timeout imposed by the evaluator is 241 seconds per candidate's test suite. This test timeout is distinct from the 180-second model-call deadline. Process limits are 30,720 MB for address space and data and 10 MB for the stack; the enclosing container imposes the stricter total memory cap of 3 GB, two CPUs and 256 processes, with a 512-MB temporary filesystem.

The initial factorial environment uses NumPy 1.26.4, pandas 2.2.3, SciPy 1.14.1, scikit-learn 1.5.2, Matplotlib 3.9.2 and NLTK 3.8.1. The control-stage additions include BeautifulSoup 4.12.3, OpenCV-headless 4.9.0.80, mechanize 0.4.10, rsa 4.9, pyquery 2.0.0 and statsmodels 0.14.1. NLTK's \texttt{punkt}, \texttt{averaged\_perceptron\_tagger}, \texttt{stopwords} and \texttt{words} resources are installed before offline testing. The separate 80-task environment also restores older library interfaces, including SciPy 1.10.1, scikit-learn 1.3.1 and Matplotlib 3.7.0, and adds TensorFlow-CPU 2.15.1. Its results are therefore reported separately.

For each task, the reference control concatenates \texttt{code\_prompt} and \texttt{canonical\_solution}. The incorrect control appends a replacement \texttt{task\_func(*args, **kwargs)} that returns \texttt{None}. A task is eligible for quality analysis only when the reference passes and this incorrect program fails. This simple negative control is insensitive on /59 in the 80-task study, which is why that task is excluded from quality inference despite successful reference execution. For eligible tasks, native pass is coded 1; native failure or native timeout is coded 0. Missing generation, an incomplete evaluator run or an invalid control does not receive an invented binary outcome.

\subsection{Calculations from paired outcomes}
For a completed single-repeat contrast, let $b$ count tasks passed only by condition A, $c$ those passed only by B, and $n$ all complete eligible pairs. The quality difference is $(b-c)/n$. With $d=b+c>0$, the exact two-sided McNemar probability is
\[
p=\min\left(1,\;2\sum_{k=0}^{\min(b,c)}{d\choose k}2^{-d}\right);
\]
when $d=0$, $p=1$. Within a family of $m$ tests, order the raw probabilities as $p_{(1)}\le\cdots\le p_{(m)}$ and calculate
\[
p^{\mathrm{Holm}}_{(i)}
=\min\left(1,\max_{j\le i}(m-j+1)p_{(j)}\right).
\]
Here $m=4$ in the initial factorial study and $m=2$ in the layout study. The reported 95\% intervals are the 2.5th and 97.5th percentiles of 10,000 resampled task means, using NumPy's generator with seed 20260908. All conditions of a sampled task remain together. The intervals are descriptive; the Holm-adjusted tests determine the planned rejections.

If $s$ successes and $u$ unavailable outcomes occur among $N$ assignments, the extreme success-rate bounds are $[s/N,(s+u)/N]$. For a paired difference, each unknown binary outcome is separately set to 0 or 1 to obtain the minimum and maximum possible difference. These bounds describe missing-outcome uncertainty, not sampling uncertainty.

\subsection{SCC adaptation and feedback}
The 2024 SCC workflow begins with an analyst's concise JSON plan, followed by the developer's Python program. Its tester then generates a \texttt{check(candidate)} function containing up to five printed example calls and is explicitly instructed not to write assertions. The controller executes the program and generated check together. Exception-free execution produces its internal success message; an exception or timeout produces feedback for the developer. Printed answers are not matched automatically against expected values. The inner execution limit is 10 seconds; the containing process has a separate 45-second limit. Neither reference solutions nor benchmark tests enter this feedback loop.

The configured limit of two developer iterations permits one initial program and at most one repair: analyst, developer, tester, and optionally developer again. The last allowed developer response ends the session without another tester call. The author's extractor first uses the first fenced code block; if absent, it attempts to recover code from function definitions. The session retains the most recent program with an identifiable top-level function, or an error result if none exists. The model receives each role's complete message history serialized as a JSON conversation. The original API controls for temperature, top-$p$ and maximum output tokens are not enforced by the subscription transport. Thus the pilot assesses this explicit software adaptation; its internal stopping message is not the separately measured BigCodeBench outcome.

\subsection{Repeated factorial study and amended SCC comparison}
\label{app:expanded-studies}
The Luna factorial allocation reuses the 200 initial task IDs and adds 800 IDs from the original randomized order, excluding development tasks, the 80 layout tasks and /0--/7. The previously studied 200 tasks are not newly held out. All model requests are fresh; historical programs are not pooled into the new study. The 15,000 assignments yield 14,961 complete candidates. Reference controls completed before generation permit quality measurement on 985 tasks. All assigned rows, including the 15 control-ineligible tasks, remain in the accounting.

\begin{table}[htbp]\centering\small
\caption{Luna allocations. A block is one task and one repeat label, with every method assigned. Completion counts refer to programs, not successful test outcomes.}
\begin{tabularx}{\textwidth}{@{}lXX@{}}\toprule
 & Factorial study & Amended SCC comparison\\\midrule
Task IDs & 1,000 & 956 from the same allocation\\
Repeats per task & Three & One to three selected labels\\
Conditions & D, SN, SR, MN, MR & Fresh SR, fresh SN, SCC\\
Assignments & 15,000 & 6,000 in 2,000 blocks\\
Maximum model calls & 27,000 & 12,000\\
Quality-test family & Four contrasts & SCC--SR and SCC--SN\\\bottomrule
\end{tabularx}\end{table}

Both studies use GPT-5.6 Luna, medium reasoning, Codex CLI 0.153.4, a 600-second deadline per model call and at most eight workers. An input guard rejects requests exceeding 65,536 UTF-8 bytes instead of shortening them. The full supplied context is retained. Submitted attempts are not retried; after an interruption, only previously untouched assignments may start. A completed malformed answer yields an observed empty candidate, while an absent completed answer remains unavailable. A model alias and repeated requests do not establish immutable weights or independent provider sampling.

For recorded input $I$, cached input $C\subseteq I$ and output $O$, Luna's API-equivalent valuations are
\[
V_L=\frac{0.20(I-C)+0.02C+1.20O}{10^6},\qquad
V_{L,0}=\frac{0.20I+1.20O}{10^6}\quad\mathrm{USD}.
\]
The standard API rates were checked on 13 September 2026~\cite{openaipricing2026}. They differ from the initial mini-study tariff. These quantities value measured subscription-channel counters at published rates; they are neither invoices nor total research costs. Output includes any separately exposed reasoning count, which is not added again. Unknown usage is not zero-imputed.

The original SCC allocation contains 3,000 randomized task-repeat blocks and 9,000 method cells. An administrative budget amendment fixes the first 2,000 blocks in that schedule before inspecting SCC quality. This prefix contains 956 tasks: 212 with one selected repeat label, 444 with two and 300 with three. Its 941 control-eligible tasks form the possible quality-analysis population. The remaining 3,000 method cells are administratively unselected, outside the 6,000-cell denominator, and retained in the original ledger. Eight previously submitted interrupted attempts remain unavailable and are not regenerated. The amendment follows the start of generation; it is not described as a fully preregistered new experiment. The selected prefix is complete for reporting: SCC has 1,600 completed workflows (396 Docker or infrastructure failures and 4 paused attempts), SR has 1,997 (3 paused), and SN has 1,998 (1 paused and 1 stream or infrastructure failure). Thus 5,595 candidates are complete, with 397 infrastructure failures and 8 paused attempts; every complete candidate received the official native evaluation. These counts refer to completed programs or workflows, not successful test outcomes.

For the factorial study, a task enters a quality contrast only if all three repeats are observed in both conditions. The difference $d_i$ is the difference of the two three-repeat means. A two-sided paired $t$-test tests the mean of these task differences. The four contrasts in Table~\ref{tab:scale-contrasts} form one fixed Holm family at $\alpha=0.05$. Each descriptive interval resamples whole tasks 10,000 times with NumPy 1.26.4's PCG64 generator, seed 20260909, preserving the frozen task order. Percentile endpoints use linear interpolation. These rules are separate from the four exact McNemar tests in the initial mini study.

The amended SCC estimator permits the unequal number of selected repeats. For a task $i$ and comparator $B$, let $K_i^B$ contain its selected repeat labels for which both SCC and $B$ have observed native outcomes under valid controls. Its contribution is
\[
d_i^B=\frac{1}{|K_i^B|}\sum_{r\in K_i^B}
\bigl(Y_{ir}(\mathrm{SCC})-Y_{ir}(B)\bigr).
\]
Only tasks with $|K_i^B|>0$ contribute, and each contributing task has equal weight. Thus a task with three selected repeats does not receive three times the weight of a task with one. SCC--SR and SCC--SN form a fixed two-test Holm family. Task-bootstrap intervals use 10,000 PCG64 draws, seed 20260911, and linear percentile endpoints. The original SCC estimator, requiring three complete repeats per condition, is retained separately with its own two-test family and original Python Mersenne-Twister bootstrap. Its original 9,000-assignment missingness bounds are not relabelled as prefix bounds. Neither analysis treats repeats as independent tasks or estimates best-of-three success.

For either task-mean estimator, the test statistic is $t=\bar d/(s_d/\sqrt n)$ when $n\geq2$ and $s_d>0$, with $n-1$ reference degrees of freedom. This is the specified large-sample approximation. An unestimable SCC test retains its place in the two-test family, using 1 for the adjustment input and supporting no rejection. A constant zero vector with at least two tasks has $p=1$. The frozen factorial analyzer encodes a nonzero constant vector with a degeneracy flag and $p=0$; this is not an ordinary finite $t$-test. No such degeneracy occurs in its four observed contrasts. No analysis introduces a quality non-inferiority margin or joint quality--cost decision.

Missingness bounds and sampling intervals answer different questions. For the SCC prefix, every task's denominator is its number of \emph{assigned} selected repeats, including unavailable outcomes. Unknown binary endpoints are completed in all favourable or all unfavourable ways, then task differences are averaged with equal task weight. Bounds are given for the 941 eligible selected tasks and for all 956 selected tasks. They describe the range allowed by missing data; they are not confidence intervals or a correction for non-random missingness. The factorial bounds likewise preserve its full assigned denominators.

Resource pairing does not require reference-control eligibility. The factorial rule requires all three workflows in both conditions to finish with complete turn counters. The amended SCC rule retains a selected repeat pair only when both entire workflows finish and every submitted turn has valid counters. Resources are summed across turns, then averaged over these matched repeats within task. Tasks again receive equal weight. Known usage from incomplete workflows is reported separately as a subtotal. The original SCC resource estimator, based on known submitted-turn totals across three repeats, remains separate and may include attempts without a completed candidate.

For paired task means $A_i$ and $B_i$, two distinct resource ratios are
\[
R_{\mathrm{pooled}}=\frac{\sum_i A_i}{\sum_i B_i},\qquad
R_{\mathrm{task}}=\frac{1}{n_R}\sum_{i:B_i\ne0}\frac{A_i}{B_i},
\]
where $n_R$ counts nonzero denominators. The factorial study uses seed 20260909 for intervals on the pooled ratio and mean paired difference. The amended SCC analysis uses seeds 20260912, 20260913 and 20260914 for the difference, mean task ratio and pooled ratio, respectively. Each statistic starts a fresh generator and retains all 10,000 draws. The original SCC analysis reports only a point estimate for the pooled ratio; its interval for the mean task ratio concerns a different quantity. Tables name the estimator and complete-pair count. Resource intervals remain descriptive and do not enter the quality-test families.

Source-group sensitivity resamples whole groups under four predefined partitions: the combined instruction/code rule and instruction-only thresholds of 0.50, 0.70 and 0.90. Group draws use seed 20260911 and preserve task weighting by dividing the summed sampled task differences by the sampled task count. No source-similarity exclusion is introduced. These intervals supplement the registered task-based analyses. SCC remains a comparison of complete workflows, including feedback from its generated checks; a role-excised SCC condition is not evaluated.
